\begin{center}
\LARGE
Cryptographic Aspects Of Real Quadratic Congruence
                        Function Fields 
\end{center}

\begin{center}
\Large
Andreas Stein
\end{center}

\noindent
Date:  July 17th (Wednesday)\\
Time: 15:00-15:30\\

\begin{center}
\large
Abstract
\end{center}

We show how the theory of real quadratic congruence function fields can be used to produce a
secure key distribution protocol. The technique is similar to that advocated by Diffie and Hellman
in 1976, but instead of making use of a group for its underlying structure, makes use of a structure
which is ``almost'' a group. The method is an extension of the recent ideas of Scheidler, Buchmann
and Williams, but, because it is implemented in these function fields, several of the difficulties
with their protocol can be eliminated. A description of the protocol is provided, together with a
discussion of the difficulty of the so-called discrete logarithm problem (DLP) in real quadratic
congruence function fields. In this context, we will point out that the DLP for real quadratic
congruence function fields of genus one is equivalent to the DLP for elliptic curves over finite
fields. 

